%% ch03 --- Typowanie: adnotacje, checkery i gdzie checker klamie
%%
%% Napisany. Opis, ktory tu stal, jest w tools/chapters.json; to, czemu
%% pisanie zaprzeczylo, jest w CLAUDE.md pod *Resolved questions*.

\chapter{Typowanie: adnotacje, checkery i gdzie checker kłamie}
\label{ch:typing}

\noindent
Spędziłeś dekadę z kompilatorem, który odmawia. Napisz \code{int x =
"8080";} i nie powstanie nic: żadnego assembly, żadnego przebiegu testów,
żadnego wdrożenia. System typów nie jest poradą, jest bramką, a każdy twój
nawyk dotyczący typów ukształtował się po jej drugiej stronie.

Python ma adnotacje i nie ma bramki. Adnotacja to twierdzenie zapisane w
kodzie źródłowym, interpreter nigdy go nie czyta, a program i tak się
wykonuje. Jeśli chcesz, żeby coś odmawiało, instalujesz to, konfigurujesz i
uruchamiasz sam \dash{} drugi kompilator, którego musisz zaprosić. Ten
rozdział jest o tym, co to zmienia, i o miejscach, w których zaproszony
checker po cichu się z tobą zgadza i nie ma racji.

Pożycza z wcześniejszych rozdziałów jedno: środowisko zdolne uruchomić
listing, czyli Rozdział~\ref{ch:packaging}. Nic więcej tutaj nie potrzebuje
Rozdziałów \ref{ch:runtime} ani~\ref{ch:packaging}, a \pkg{pydantic}
pojawia się tylko na tyle, na ile potrzebuje go argument tego rozdziału
\dash{} głębia jest w Rozdziałach~\ref{ch:services}
i~\ref{ch:aitoolkit}.

\section{Adnotacja to twierdzenie}
\label{sec:ch03-claim}
\index{typowanie!adnotacja jako twierdzenie}

Przeczytaj poniższy listing, zanim go uruchomisz, i zdecyduj, czego się
spodziewasz. Trzy wywołania, każde łamiące adnotację nad sobą. Ile z nich
rzuci wyjątek?

\pyfile{code/ch03/hint_is_a_claim.py}{Trzy złamane adnotacje i to, co
interpreter ma o nich do powiedzenia}{lst:ch03-claim}

\noindent
Żadne. Oto, co wypisuje:

\transcript{ch03-hint-is-a-claim}

\noindent
\api{shout} jest zaadnotowane tak, by przyjmować \code{str}, i dostaje
\code{int}; mimo to zwraca \code{str}, bo f-string sformatuje cokolwiek mu
podasz. \api{port\_of} jest zaadnotowane tak, by zwracać \code{int}, a
zwraca \code{str}, i wywołujący, który doda do tego tysiąc, dowie się o tym
zupełnie gdzie indziej. A same adnotacje to zwykłe dane wiszące przy
obiekcie funkcji, czytelne i zapisywalne w czasie wykonania \dash{} tak
\pkg{pydantic} i \pkg{FastAPI} czytają twoje sygnatury \dash{} co znaczy
też, że adnotację można po fakcie zastąpić kłamstwem i nic tego nie
zauważy.

\begin{csbox}
Najbliższym odpowiednikiem w C\# nie jest \code{dynamic}. Jest nim baza
kodu, w której każda nazwa typu jest komentarzem: nazwy są na miejscu, przez
większość czasu są prawdziwe, a środowisko uruchomieniowe nigdy na żadną nie
spojrzało. Nullable reference types to najbliższa rzecz, której już
używałeś \dash{} adnotacja, o której rozumuje kompilator, a CLR nie.
\end{csbox}

\begin{trapbox}
\textbf{\code{str} znaczy, że wartość jest napisem.} Nie znaczy. Znaczy, że
ktoś miał na myśli napis. Parametr zaadnotowany \code{str} może w czasie
wykonania być \code{None}, może być \code{int}, może być otwartym gniazdem, i
nic w języku tego nie powie. Wszystko, co daje ci adnotacja, pochodzi z
narzędzia, które sam uruchomiłeś, a następna sekcja jest o zaproszeniu go.
\end{trapbox}

\mermaidfig{ch03-two-compilers}{Twierdzenie, uruchomienie i drugi
kompilator, którego trzeba zaprosić}{adnotacja, interpreter i checker}

\section{Mapa}
\label{sec:ch03-map}
\index{typowanie!mapa z C\# na Pythona}

Większość słownika przenosi się wprost i poniższa tabela jest tą częścią,
którą czyta się raz i zatrzymuje. Dwa wiersze się nie przenoszą i dostają
resztę tej sekcji.

\begin{center}
\small
\begin{tabularx}{\linewidth}{@{}lX@{}}
\toprule
\textbf{C\#} & \textbf{Python} \\
\midrule
\code{int}, \code{string}, \code{bool} & \code{int}, \code{str}, \code{bool}
  \\
\code{List<int>}, \code{Dictionary<string, int>} & \code{list[int]},
  \code{dict[str, int]} \\
\code{string?}, nullable reference types & \code{str | None}, prawdziwa unia,
  którą trzeba zawęzić \\
\code{IEnumerable<T>}, \code{IReadOnlyList<T>} & \api{Iterable[T]},
  \api{Sequence[T]}, z \pkg{collections.abc} \\
Interfejs, implementowany przez deklarację & \api{Protocol}, spełniany
  kształtem \\
\code{record}, DTO, POCO & \code{dataclass}, \api{TypedDict},
  \api{NamedTuple}, model \pkg{pydantic} \\
\code{enum} & \api{Enum} albo \api{Literal} dla zamkniętego zbioru wartości
  \\
Wzorzec \code{is}, \code{out var} & \code{isinstance} albo własny predykat
  \api{TypeIs} \\
\code{object} & \code{object}, który trzeba zawęzić przed użyciem \\
\code{dynamic} & \api{Any}, co nie jest tym samym \\
\code{where T : IComparable} & ograniczenie na parametrze typowym \\
\code{in T}, \code{out T} & nic: wariancja jest wnioskowana \\
\bottomrule
\end{tabularx}
\end{center}

\noindent
Generyki to pierwszy wiersz wart zatrzymania. Współczesny zapis deklaruje
parametr na funkcji albo klasie, bez \api{TypeVar} i bez importu:

\pyregion{code/ch03/the_map.py}{generics}{Dwie klasy generyczne i
wariancja, której nikt nie zadeklarował}{lst:ch03-generics}

\noindent
Żadna z klas nie napisała \code{out} ani \code{in}. Checker czyta, w których
pozycjach pojawia się parametr \dash{} \api{Reader} tylko zwraca \code{T},
\api{Sink} tylko przyjmuje \dash{} i wnioskuje z tego kowariancję oraz
kontrawariancję. Oba wywołania w \api{variance} są przyjęte, a zamiana
argumentów miejscami jest odrzucana. To ta sama idea co słowa kluczowe w
C\#, z usuniętą deklaracją, a rekompensatą jest to, że nie da się już
pomylić adnotacji.

\begin{trapbox}
\textbf{\code{List[int]} i \code{Optional[str]}.} To zapis mniej więcej z
roku 2019, to właśnie da ci wynik wyszukiwania, a wbudowane generyki
zastąpiły go lata temu. Pisz \code{list[int]} i \code{str | None}. Jeśli
czytasz tutorial, który importuje \api{List} z \pkg{typing}, załóż, że reszta
strony jest w tym samym wieku.
\end{trapbox}

Zawężanie to drugi wiersz. \code{str | None} jest unią i checker nie pozwoli
ci o tym zapomnieć, więc wartość jest naprawdę dwoma typami, dopóki gałąź
ich nie rozdzieli \dash{} co jest ostrzejsze niż nullable reference types i
znacznie łatwiejsze w życiu. Tam, gdzie zwykły \code{isinstance} nie
wystarczy, możesz napisać własny predykat, a zapisy są dwa i różnią się
jedną ważną rzeczą:

\pyregion{code/ch03/the_map.py}{narrowing}{\api{TypeIs} zawęża obie gałęzie;
\api{TypeGuard} zawęża jedną}{lst:ch03-narrowing}

\noindent
\api{TypeGuard} mówi checkerowi, co jest prawdą, gdy predykat zwróci prawdę,
i nic o drugiej gałęzi, więc drugie sprawdzenie w
\api{length\_with\_typeguard} nie jest ozdobą: usuń je, a linia przestanie
przechodzić kontrolę typów. \api{TypeIs} zawęża obie gałęzie, co niemal
zawsze jest tym, co miałeś na myśli.

\begin{note}
\api{TypeIs} kupuje to, żądając czegoś, czego \api{TypeGuard} nie żąda: typ,
do którego zawężasz, musi być przypisywalny do typu, \emph{od} którego
zawężasz. Dlatego listing przyjmuje \api{Sequence}, a nie \code{list}.
\code{list} jest niezmienniczy, więc \code{list[str]} nie jest
\code{list[object]}, a \code{list[object] -> TypeIs[list[str]]} zostaje
odrzucone tam, gdzie para z \api{Sequence} przechodzi. Wariancja trzy
akapity wyżej i zawężanie tutaj to ta sama reguła widziana dwa razy.
\end{note}

\begin{exercise}{e03_01_narrow}{Zawęź typ, a nie tylko odpowiedz na pytanie}
Napisz predykat odpowiadający, czy każdy element sekwencji jest typu
\code{str}. Test sprawdza odpowiedzi \emph{oraz} adnotację zwracaną, bo
predykat zwracający \code{bool} nie mówi checkerowi nic: odpowiedź w czasie
wykonania i statyczne zawężenie to dwie osobne prace, a to ćwiczenie chce
obu. Sygnatura w pliku startowym jest celowo tą niewłaściwą.
\end{exercise}

\section{To samo DTO na cztery sposoby}
\label{sec:ch03-dto}
\index{typowanie!dataclass, TypedDict, NamedTuple}

Rekord, DTO i POCO to w C\# jedna decyzja. W Pythonie są cztery, a wybór nie
jest kwestią gustu. Zanim przeczytasz listing: ten sam JSON trafia do
wszystkich czterech, a port jest w nim napisem \code{"8080"}. Zapisz, która z
czwórki twoim zdaniem go odrzuci.

\pyfile{code/ch03/dto_four_ways.py}{Jeden kształt, cztery sposoby na jego zadeklarowanie}{lst:ch03-dto}

\transcript{ch03-dto-four-ways}

\noindent
Jedna z czterech. \code{@dataclass} generuje \api{\_\_init\_\_},
\api{\_\_eq\_\_} i \api{\_\_repr\_\_}, a nie waliduje niczego;
\api{TypedDict} to kształt słownika, który już istnieje, więc w czasie
wykonania nie ma tam klasy ani instancji, tylko \code{dict};
\api{NamedTuple} jest prawdziwą krotką, dlatego porównuje się równo z
anonimową. Tylko model \pkg{pydantic} patrzy na wartość i tylko on z całej
czwórki robi przy tworzeniu cokolwiek poza przypisaniem pól.

Zwróć uwagę, co checker zobaczył, a czego nie. JSON przychodzi jako
\api{Any}, więc żadna z tych czterech konstrukcji nie potrzebowała
komentarza, żeby przejść przez \api{pyright}. Napisz tę samą błędną wartość
jako literał, dwie linie niżej, a zostanie złapana w obu liniach. Nic w
błędzie się nie zmieniło; zmieniło się tylko to, skąd przyszła wartość.

\begin{trapbox}
\textbf{\code{@dataclass} waliduje swoje pola, jak rekord z konstruktorem.}
Nie waliduje, i to jest najdroższe założenie, jakie inżynier C\# przynosi do
Pythona. Dataclass to generator boilerplate'u, nie bramka. Jeśli wartości
przychodzą spoza procesu, sam dataclass jest złym narzędziem, a dobrym jest
\pkg{pydantic} \dash{} o czym jest cała \S\ref{sec:ch03-boundary}.
\end{trapbox}

\begin{center}
\small
\begin{tabularx}{\linewidth}{@{}lXl@{}}
\toprule
\textbf{Kształt} & \textbf{Sięgnij po niego, gdy} & \textbf{Sprawdza?} \\
\midrule
\code{dataclass} & wartość jest już twoja i już poprawna & nie \\
\api{TypedDict} & \code{dict} istnieje, a ty chcesz zapisać jego kształt &
  nie \\
\api{NamedTuple} & chcesz taniej krotki z nazwami & nie \\
\pkg{pydantic} & wartość przyszła spoza tego procesu & tak \\
\bottomrule
\end{tabularx}
\end{center}

\section{Protocol jest spełniany kształtem}
\label{sec:ch03-protocol}
\index{typowanie!Protocol i typowanie strukturalne}

Interfejs w C\# to kontrakt podpisany przez implementującego: klasa spełnia
\code{IDisposable}, bo tak powiedziała. \api{Protocol} to kontrakt pisany
przez \emph{konsumenta}, a klasa spełnia go, mając odpowiednie składowe.

\pyfile{code/ch03/structural.py}{Interfejs, o którym implementujący nigdy nie słyszeli}{lst:ch03-protocol}

\noindent
\api{Connection} nie importuje \api{Closable} i nie robi tego również
\api{StringIO}, który trafił do interpretera lata przed tym, jak napisałeś
swój protokół. Oba przechodzą przez \api{close\_all}. To właśnie kupuje
typowanie strukturalne, a C\# nie potrafi: możesz napisać interfejs dla typu,
którego nie jesteś właścicielem, po fakcie i bez adaptera.

Połowa działająca w czasie wykonania jest słabsza od statycznej, a wiedza o
tym, jak dokładnie słabsza, jest warta więcej niż sama funkcja.
\api{runtime\_checkable} sprawia, że \code{isinstance} działa wobec
protokołu, a \code{isinstance} pyta wyłącznie o to, czy \emph{atrybuty są
obecne}. Nie patrzy na sygnatury. Więc \api{Detonator}, którego
\api{close} bierze parametr i przez to w ogóle protokołu nie spełnia,
przechodzi sprawdzenie w czasie wykonania \dash{} a \api{pyright} nie
przepuszcza tej linii po cichu, dlatego listing niesie komentarz
wyciszający, który nazywa regułę. Używaj \code{isinstance} wobec protokołu
jako zgrubnego filtra; po odpowiedź idź do checkera.

\begin{exercise}{e03_02_protocol}{Interfejs, o którym implementujący nigdy nie słyszeli}
Napisz protokół \api{Closable} sprawdzalny w czasie wykonania i funkcję,
która zamyka każdy element go spełniający, resztę zostawia w spokoju i
zwraca, ile zamknęła. Nic, co poda test, nie będzie dziedziczyć po twoim
protokole, włącznie z jedną klasą z biblioteki standardowej. I o to chodzi.
\end{exercise}

\section{Gdzie checker kłamie}
\label{sec:ch03-lies}
\index{typowanie!propagacja Any}

Zaproś checkera, a złapie bardzo dużo. Powie ci też, w czterech częstych
sytuacjach, że plik jest czysty, kiedy nie jest. Poniższy listing ma
\val{ch03.lies.functions} funkcje. Każda deklaruje, że zwraca \code{int}.
Każda zwraca \code{str}.

\pyfile{code/ch03/where_it_lies.py}{Cztery sposoby, żeby checker przestał widzieć}{lst:ch03-lies}

\noindent
Pierwsze dwa to \api{Any} przychodzące bez niczyjej decyzji.
\api{json.loads} jest zaadnotowane tak, by zwracać \api{Any}, więc wszystko
poniżej jest \api{Any}, a \api{Any} spełnia dowolny zadeklarowany typ
\dash{} to znaczy to słowo w systemie typów, a nie błąd narzędzia. Pakiet,
który nie dostarcza informacji o typach, robi to samo ze wszystkim, co ci
poda. Ostatnie dwa to ty mówiący checkerowi, żeby przestał patrzeć:
\api{cast} nie generuje kodu i niczego nie sprawdza, a \code{type: ignore}
usuwa błąd z raportu, nie usuwając go z programu. Oba są uprawnione i oba
przenoszą odpowiedzialność na ciebie.

Co więc dwa checkery przypięte przez tę książkę robią z tym plikiem? To jest
generowane przez \code{code/measure/checkers.py}, który uruchamia oba i
zapisuje, co powiedział każdy z nich:

\transcript{ch03-checkers}

\noindent
\api{pyright} w trybie strict raportuje \val{ch03.pyright.errors}.
\api{mypy} pod \code{-{}-strict} raportuje \val{ch03.mypy.errors} \dash{}
oba razy chodzi o te dwa zwroty \api{Any}, a żaden z nich nie dotyczy
\api{cast} ani \code{type: ignore}, które w obu narzędziach milczą z
założenia. Ta różnica nie jest wadą żadnego z nich. \api{pyright} stosuje
system typów, w którym \api{Any} jest zgodne ze wszystkim; profil strict
\api{mypy} dokłada opinię, której system typów nie ma: że funkcja
deklarująca \code{int} nie powinna po cichu zwracać \api{Any}.

\begin{trapbox}
\textbf{\api{Any} jest jak \code{dynamic}: coś, co potrafię ogrodzić.}
\api{Any} się propaguje. Jedna taka wartość zamienia wszystko, co z niej
policzone, również w \api{Any}, po cichu, na całej drodze, a raport przez ten
czas pozostaje czysty. \code{dynamic} w C\# przynajmniej wywraca się głośno w
miejscu użycia. \api{Any} nie wywraca się wcale; zatrzymuje sprawdzanie, a
program leci dalej, myląc się gdzieś później.
\end{trapbox}

A checker nie jest jedynym drugim kompilatorem, którego możesz zaprosić. Oto
funkcja, której adnotacja jest poprawna, której checker jest zadowolony i
która jest błędna:

\pyfile{code/ch03/defaults.py}{Argument domyślny, wyliczany raz}{lst:ch03-defaults}

\transcript{ch03-defaults}

\begin{trapbox}
\textbf{\code{def f(items: list[str] = [])} to opcjonalny parametr listowy.}
Wartość domyślna jest wyliczana raz, gdy wykonuje się linia \code{def}, i ta
jedna lista jest współdzielona przez każde wywołanie, które nie poda
własnej. Adnotacja jest poprawna, a funkcja błędna. Podawaj \code{None} i
buduj listę w środku.
\end{trapbox}

\noindent
Żaden z przypiętych checkerów nie mówi o tym pliku nic. Linter mówi:

\transcript{ch03-three-tools}

\noindent
\api{pyright} znajduje \val{ch03.defaults.pyright}, \api{mypy} znajduje
\val{ch03.defaults.mypy}, a \pkg{ruff} znajduje \val{ch03.defaults.ruff}.
Wniosek wychodzi poza tę jedną regułę: checker typów odpowiada na pytania o
typy, a spora część tego, co psuje się w Pythonie, nie jest pytaniem o typy.
Uruchamiaj oba.

\begin{projectbox}
Ten rozdział nie dokłada etapu do \code{trace-assert} \dash{} pierwszy jest
w Rozdziale~\ref{ch:testing} \dash{} ale projekt płaci rachunek tego
rozdziału od pierwszego commita. Otwórz
\code{code/src/trace\_assert/\_\_init\_\_.py} i przeczytaj, dlaczego payload
w \api{Event} używa nazwanej fabryki, a nie \code{default\_factory=dict}. Pod
ścisłym checkerem goły \code{dict} jest typu \code{dict[Unknown, Unknown]}, a
naprawą jest funkcja z typem zwracanym. Ta sama skarga padła znowu przy
pisaniu \code{measure/checkers.py} i została załatwiona jednym
udokumentowanym \api{cast} \dash{} jedną z czterech dziur wyżej, użytą
celowo i w jednym miejscu.
\end{projectbox}

\begin{exercise}{e03_03_defaults}{Domyślna wartość, która nie jest tym, na co wygląda}
To ćwiczenie zaczyna się błędne, a nie niedokończone: funkcja jest
zaadnotowana dokładnie tak, jak zaadnotowałbyś ją w C\#, checker jest z niej
zadowolony, a test wywraca się na drugim wywołaniu. Napraw ją tak, żeby dwaj
wywołujący, którzy nic nie podadzą, nie dzielili jednej listy.
\end{exercise}

\section{Który checker i konfiguracja do skopiowania}
\label{sec:ch03-config}
\index{typowanie!pyright i mypy}

Te dwa nie są wymienne i poprzednia sekcja mówi dlaczego. \api{pyright}
prawdopodobnie już działa w twoim edytorze, a jego tryb strict jest naprawdę
wysoko postawioną poprzeczką. \api{mypy} jest starszy z tej dwójki o ponad
dekadę, a jego profil strict niesie garść opinii, których \api{pyright} nie
ma, i z których \code{no-any-return} jest tą zmierzoną w tym rozdziale.
Który z nich jest szybszy na twoim drzewie, tego ta książka nie zmierzyła,
więc nie przyjmuj twierdzenia o szybkości od nikogo, kto też tego nie
zrobił.

Ta książka przypina oba i na każdym pushu puszcza \api{pyright} po całym
drzewie. Rekomendacja, która wynika z pomiaru, a nie z gustu: uruchamiaj w CI
jedno z nich po całym drzewie, a jeśli tym jednym jest \api{pyright}, wiedz,
że \api{Any} wyciekające z zadeklarowanego typu zwracanego to klasa defektów,
której nie zgłosi. Zacznij stąd, w \code{pyproject.toml}:

\begin{tomlcode}
[tool.pyright]
typeCheckingMode = "strict"
include = ["src", "tests"]

[tool.mypy]
strict = true
files = ["src"]
\end{tomlcode}

\noindent
Oba narzędzia biorą wersję interpretera ze środowiska, w którym działają,
więc przypnij ją obok tych bloków, jeśli twój projekt potrzebuje innej;
\code{code/pyproject.toml} tej książki jest wersją tego wszystkiego do
skopiowania. Dwie uwagi o życiu z trybem strict, obie wyniesione z kodu tej
książki. Włączaj go na starcie projektu albo pakiet po pakiecie; włączenie go
naraz w istniejącej bazie kodu daje liczbę diagnostyk, której nikt nie
przeczyta. A kiedy tryb strict skarży się, że typ jest \emph{częściowo
nieznany}, prawie nigdy nie jest drobiazgowy \dash{} znalazł miejsce, w
którym wydaje ci się, że masz typ, a nie masz.

\section{Waliduj na granicy, ufaj w środku}
\label{sec:ch03-boundary}
\index{typowanie!walidacja na granicy}

Wszystko powyżej zwija się do jednej reguły i jest to jedyna rzecz w tym
rozdziale warta zapamiętania.

\mermaidfig{ch03-boundary}{Jedno miejsce waliduje; wszędzie po nim adnotacje
znaczą to, co mówią}{waliduj na granicy}

\noindent
Wewnątrz własnego kodu adnotacja jest warta bardzo dużo: checker ją czyta,
edytor z niej podpowiada, a błędną łapie w miejscu, w którym ją piszesz. Na
krawędziach \dash{} JSON, środowisko, formularz, wiadomość z kolejki,
odpowiedź serwisu \dash{} adnotacja nie jest warta nic, bo wartość
przychodzi jako \api{Any}, a checker przestał patrzeć linię wcześniej.

Umieść więc sprawdzanie tam, gdzie wartości wchodzą, raz, w miejscu, które
umiesz nazwać, a wszystko po nim niech będzie zwykłym typowanym Pythonem. Do
tego służy \pkg{pydantic} i dlatego jest na każdej liście zależności w tej
dziedzinie, i dlatego tabela w \S\ref{sec:ch03-dto} ma kolumnę
odpowiadającą \emph{nie} trzy razy: pozostałe trzy kształty są dla wartości,
które już są twoje. Rozdziały~\ref{ch:services} i~\ref{ch:aitoolkit} biorą to
na poważnie; temu rozdziałowi wystarczy, że wiesz, gdzie przebiega linia.

\begin{csbox}
Stosujesz tę regułę od lat, nie umieszczając jej nigdzie, bo ASP.NET Core
umieszcza ją za ciebie: model binding jest granicą, a związany model jest
zwalidowany, zanim zobaczy go twoja metoda akcji. Nic w Pythonie nie stawia
granicy nigdzie domyślnie. Reguła jest ta sama; nowe jest to, że jej miejsce
jest teraz twoją decyzją, a baza kodu, która nigdy tej decyzji nie podjęła,
ma granice rozsiane po całości.
\end{csbox}

\begin{exercise}{e03_04_boundary}{Waliduj na granicy}
Zamień dokument JSON w typowany obiekt albo go odrzuć. Test podaje twojej
funkcji taki JSON, jaki naprawdę przychodzi z granicy \dash{} port będący
napisem, port niebędący w ogóle liczbą, dokument bez klucza \dash{} i
sprawdza, że to, co wraca, ma typy, do których się przyznaje, albo że nie
wraca nic. Napisz to ręcznie albo z \pkg{pydantic}; testowi zależy na
zachowaniu, a nie na drodze.
\end{exercise}

\section{Co wynieść z tego rozdziału}
\label{sec:ch03-summary}

Adnotacja to twierdzenie, którego środowisko uruchomieniowe nigdy nie
sprawdza, a checker to drugi kompilator, którego trzeba zaprosić,
skonfigurować i uruchomić. Raz zaproszony zarabia na siebie: mapa wyżej
przenosi się niemal w całości, \api{Protocol} robi rzecz, której interfejsy
C\# nie potrafią, a tryb strict znajduje prawdziwe defekty. Ale przestaje
patrzeć na \api{Any}, a dwa checkery, które możesz uruchomić, nie przestają w
tym samym miejscu. Umieść walidację na granicy, trzymaj wnętrze typowane, a
każdy \api{cast} i każdy \code{type: ignore} traktuj jak obietnicę, którą
właśnie złożyłeś samemu sobie.
